\section{Conclusion}

We use Double DQN to divide a fixed budget between graph repair and rule
acquisition. Deletion and augmentation prompts generate rule candidates.
The tests follow three stages: scheduling, generation, and rule execution.

The count penalty charges for temporary duplicate edges and blocks useful
relation repair. The rate penalty raises Double-DQN triple F1 from 94.67\%
to 99.49\% on new corruption scenarios. Acquire-then-deficit reaches 99.91\%.
Removing the action mask lowers F1 by 2.50 points after correction for
multiple tests. Small-count DDQN reaches 99.05\% F1 at the base size; rate DDQN reaches
99.52\%. Their paired difference is $-0.47$ points
(95\% CI: $[-1.73,0.46]$). These tests show how reward size
and allowed actions affect repair.

The two prompts produce different candidates on the same documents.
Augmentation gives more constraints at equal call budgets.
Direct execution removes ten designed defects. Deletion alone covers all
ten removals. The second document-development round reaches 95.69\% F1,
removes seven injected items, and loses zero reference facts. All seven
removals use deletion type rules. The best allowed schedule has the same
final F1 as acquire-then-repair.

A new test on twenty other development documents activates source rules
and reaches 96.63\% F1. It removes 17 injected items and loses four reference
facts. Deletion alone gives the same final graphs. Requiring approval
with the current empty review sources keeps all reference facts and makes
zero repairs. The next step is to improve rule review and useful rule
coverage, then test learned scheduling against the same simple baselines.
Independent rule and edit review can check meaning. New document domains
can test transfer.
